\documentclass[11pt]{article}
\usepackage{amsmath,amssymb}
\usepackage[margin=1in]{geometry}
\usepackage{booktabs}
\usepackage{hyperref}

\title{Polarization in NeuralDMD: what we tried, and where we landed}
\author{Sgr\,A$^\ast$ dynamics / NeuralDMD polarization effort}
\date{\today}

\begin{document}
\maketitle

NeuralDMD began as a Stokes-$I$ imager. Adding linear polarization means choosing
how to represent $Q$ and $U$. This note records the four parameterizations we
tried and why we kept the last one. The test throughout is a synthetic thick ring
with a $20\%$ \emph{spiral} EVPA---an $m\!=\!2$ (four-lobe) pattern in $Q,U$---
observed with a realistic EHT array.

A good parameterization has to do three things at once:
\begin{enumerate}
  \item keep $P\equiv\sqrt{Q^2+U^2}\le I$ (a source cannot be $>100\%$ polarized);
  \item keep polarization only where there is total intensity ($Q,U\to0$ as $I\to0$);
  \item still represent the $m\!=\!2$ EVPA the data contain.
\end{enumerate}
Each choice below trades these off differently. Only the last gets all three.

\section*{Choice 1 --- Fractional $(I,\,m_\ell,\,\xi)$}
Write $Q=m_\ell I\cos2\xi$, $U=m_\ell I\sin2\xi$, with a sigmoid on the fraction
$m_\ell$ and a unit vector for the EVPA direction $(\cos2\xi,\sin2\xi)$. Both
constraints hold for free. The problem is the angle. Polarization is a spin-2
field, so the direction advances by $4\pi$ as $\xi$ turns by $2\pi$. Any ring
EVPA that rotates once around the ring---radial, azimuthal, or spiral---makes that
direction \emph{wind twice}, and a unit vector that winds must vanish somewhere,
which is a singularity gradient descent will not cross. In practice every fit
collapses to a nearly uniform EVPA ($\sim\!40^\circ$ error), regardless of network
size or training budget. It is a representation limit, not under-training. We drop
this chart.

\section*{Choice 2 --- Free direct $(I,\,Q,\,U)$}
Drop the angle and make $Q,U$ independent signed fields. The winding is gone and
the cross-hand $\chi^2$ reaches the noise floor. But nothing now ties
polarization to $I$, and the fit abuses that: it dumps ${\sim}1.3$~Jy of polarized
``haze'' off the source (truth ${\sim}0.1$) to fit the cross hands, and leaves the
ring almost unpolarized. We need a constraint back.

\section*{Choice 3 --- Free direct, plus a support penalty}
Penalize polarized flux where $I$ is faint:
$\mathcal{R}=\big\langle\sum_x P\,e^{-I/\tau}\big\rangle$. ($I$ enters through a
stop-gradient---otherwise the model cheats by brightening $I$ to switch the
penalty off, which wrecked the Stokes-$I$ image on the first try.) It does not
work. The haze and the true $m\!=\!2$ sit in different basins: to shed the haze
the fit would first have to pay the $\chi^2$ of a hazeless-but-uniform image
($\approx\!8$, against $0.4$ at the floor), which no sane penalty weight offsets,
so the optimizer stays put. A heavier penalty just erases all polarization. The
lesson: don't fight the haze with a penalty---use a parameterization where it
cannot form.

\section*{Choice 4 --- I-scaled direct $Q=I\tanh q$ \; (adopted)}
The haze exists only because free $Q,U$ are untied from $I$. So tie them, but keep
them direct:
\begin{equation*}
  Q = I\,\tanh q,\qquad U = I\,\tanh u,
\end{equation*}
with $q,u$ free signed fields. This fixes both failure modes by construction:
\begin{itemize}
  \item $Q,U\propto I$, so they vanish where $I$ vanishes---\textbf{no haze is
  possible}, and constraint~(2) holds exactly.
  \item $q,u$ are free fields, not an angle, so a four-lobe is as easy as the
  ring---\textbf{no winding}, and $m\!=\!2$ is representable.
  \item $\tanh$ caps $|q|,|u|\le1$, so $P\le\sqrt2\,I$ automatically; the last step
  to $P\le I$ is a light soft penalty.
\end{itemize}
It needs no support or compactness prior. This is the same shape KINE uses
(polarization $=I\times$ a field), but without KINE's frequency-1 smoothing, so it
keeps the bandwidth to reach $m\!=\!2$. We verified the structural guarantees
numerically; the reconstruction itself is being validated now.

\section*{Is the data even right?}
Worth ruling out, since the EVPA kept coming out wrong. Two checks say yes. First,
the truth pushed through the exact forward operator matches the measured
cross-hand data at the noise floor---so generation is correct. Second, swapping
the truth's spiral for the best-fit \emph{uniform} EVPA makes the cross-hand
$\chi^2$ about $20\times$ worse. The data clearly contain and demand the spiral;
the whole difficulty was the haze escape route, which Choice~4 closes.

\section*{Summary}
\begin{center}
\begin{tabular}{@{}lcccl@{}}
\toprule
Parameterization & $P\le I$ & no haze & reaches $m\!=\!2$ & verdict \\
\midrule
Fractional $(m_\ell,\xi)$      & free      & free              & \textbf{no} (winding) & dropped \\
Free direct $(Q,U)$            & soft      & \textbf{no} (haze)& yes                   & not enough \\
\; + support penalty           & soft      & penalty (fails)   & yes                   & barrier not cleared \\
\textbf{I-scaled} $Q=I\tanh q$ & near-free & \textbf{yes}      & \textbf{yes}          & \textbf{adopted} \\
\bottomrule
\end{tabular}
\end{center}
\noindent
Fractional has the constraints but not the bandwidth; free direct has the
bandwidth but not the constraints; the penalty can't buy the constraints back.
Tying $Q,U$ to $I$ restores the support constraint structurally while keeping the
direct chart's freedom---the only choice that meets all three requirements.
(Throughout, we never divide by $I$: $Q,U$ are always $I$ times something, and the
likelihood is on the visibilities $RR,LL,RL,LR$.)
\end{document}
